% ============================================================
%  sec6_testing.tex —— Ch6 Performance Testing and Monitoring
%  负责人：Li Hengchen (E)
%  粘贴到 Overleaf 的 sections/sec6_testing.tex
%
%  写作约定（《02_Overleaf 指南》）：
%   · 本节用 \section 起头（main.tex 不再包一层）
%   · 只写正文，不要 \documentclass / \begin{document}
%   · 数字必须来自 results/*.csv，禁止手抄、禁止估
%   · 图：文件放 report/figures/，LaTeX 里仍写 figures/xxx（相对 report/main.tex 解析）
%   · 表：tables/tab_e_loadtest.tex（由 collect_results.py --latex 生成）
%   · 特殊字符：\% $\geq$ $\times$，下划线用 \_
%
%  待 A 审核的跨人引用：
%   · \ref{sec:fault} 指向 D 的 Ch4（sec4_fault_tolerance.tex）
%   · \cite{aws2023security} 已在 references.bib；\cite{aws2024cloudwatch}
%     需 F 追加（条目见 docs/e_bib_entries.bib）
% ============================================================

\section{Performance Testing and Monitoring}\label{sec:testing}

Functional correctness alone does not establish that an onboarding pipeline is
usable: a KYC service must also absorb a traffic burst without dropping valid
applications and must isolate malformed uploads without stalling the queue. This
section defines the metrics that were instrumented, describes the fault-injected
workload and the fixture suite behind it, explains how alarms convert those
measurements into an operational response, and states the limits of what the
measurements can support.

\subsection{Observability Metrics and Instrumentation}\label{sec:metrics}

A single CloudWatch dashboard (\texttt{vericloud-dashboard}) consolidates four
families of metrics; the choice follows the operational-visibility guidance of
the AWS Well-Architected Framework~\cite{aws2023security}, while the alarm
thresholds are derived from the failure taxonomy in
Section~\ref{sec:fault}~\cite{aws2024cloudwatch}:

\begin{itemize}
  \item \textbf{Invocations and concurrency} per worker Lambda, to show that
        concurrency scales with the burst and returns to zero once the queue
        drains;
  \item \textbf{Duration} reported at p90 and p99 rather than the mean, because
        the mean hides the cold-start tail that users actually perceive;
  \item \textbf{Errors and throttles} across Lambda, DynamoDB and Step
        Functions, which distinguish application faults from capacity limits;
  \item \textbf{Queue depth} (\texttt{ApproximateNumberOfMessagesVisible}) for
        both the ingestion queue and the dead-letter queue, the leading
        indicator of back-pressure.
\end{itemize}

Latency is reported under three explicit clocks, because a single number hides
where the time is spent. \emph{Ingress latency} runs from the completion of the
client upload (\texttt{created\_at}, the S3 \texttt{eventTime}) to the moment
the pre-check Lambda starts (\texttt{received\_at}), and therefore covers event
delivery, queueing and the cold start of the first function. \emph{Service
latency} runs from \texttt{received\_at} to \texttt{processed\_at} and covers
only pipeline work. \emph{End-to-end latency} runs from the client-side
completion of the selfie upload (\texttt{t\_submit\_utc}) to
\texttt{processed\_at}, and is the number a user would experience. Separating
the three prevents the common mistake of quoting a service-internal figure as the
user-visible response time.

Percentiles are preferred to the arithmetic mean for a second reason: the
distribution of end-to-end latency in an event-driven pipeline is right-skewed,
because a small number of requests either hit a cold start or wait behind a
queue. A mean would average those cases into an unremarkable number, whereas p90
and p99 keep them visible.

\subsection{Workload Design and Fault Injection}\label{sec:workload}

The load generator dispatches 100 application packages within a 60-second
window, roughly 1.7 requests per second. In the default \emph{uniform} profile
packages are paced evenly, which measures sustained throughput without allowing
a queue to build up; a separate \emph{burst} profile releases the same 100
packages within ten seconds to observe concurrency ramp-up and cold-start
behaviour. Each package consists of one identity-document image and one selfie,
uploaded through 5-minute presigned URLs in that order; the selfie upload is the
event that triggers the workflow.

Twenty percent of the packages are deliberately corrupted (zero-byte files,
non-image payloads renamed to \texttt{.jpg}, and truncated JPEG streams), mixed
uniformly into the whole submission window rather than front-loaded, so that the
dead-letter path is exercised concurrently with the happy path. A separate
functional suite of 27 fixtures covers all four verdicts---\texttt{PASS},
\texttt{NEEDS\_HUMAN\_REVIEW}, \texttt{FAIL} and \texttt{DLQ}---as
summarised in Table~\ref{tab:fixtures}.

\begin{table}[htbp]
  \centering
  \small
  \caption{Functional fixture coverage by expected verdict.}
  \label{tab:fixtures}
  \begin{tabular}{llc}
    \toprule
    Category & Expected outcome & Cases \\
    \midrule
    Normal (same person, valid document) & \texttt{PASS} & 10 \\
    Identity mismatch & \texttt{NEEDS\_HUMAN\_REVIEW} & 3 \\
    Blurred selfie & \texttt{NEEDS\_HUMAN\_REVIEW} & 2 \\
    Face-less selfie & \texttt{NEEDS\_HUMAN\_REVIEW} & 2 \\
    Expired document & \texttt{FAIL} & 2 \\
    Missing OCR field & \texttt{FAIL} & 2 \\
    Zero-byte upload & DLQ & 2 \\
    Wrong format / truncated & DLQ & 4 \\
    \midrule
    Total & & 27 \\
    \bottomrule
  \end{tabular}
\end{table}

The two blurred selfies deserve a note: a blurred capture may still expose a face
that simply fails the similarity threshold, or it may expose no usable face at
all. Both outcomes are correct for this category and both are classified as
\texttt{NEEDS\_HUMAN\_REVIEW}, so the expected reason for these two cases is
recorded as either \texttt{face\_error} or \texttt{similarity\_below\_80}.

The fixtures are generated locally and exchanged as a content-addressed bundle:
image files are kept out of the repository for privacy, and a SHA-256 manifest
allows every team member to confirm that they are testing against byte-identical
inputs. Without that check, three copies generated independently on three
machines would silently disagree, and a defect reported by one member would not
be reproducible by another.

All experiments run in \texttt{us-east-1} on AWS Academy Learner Lab. The four
Lambdas use the Python 3.12 runtime with 1024\,MB of memory and a 60\,s timeout;
the state machine is a Standard workflow; the ingestion queue uses a 360\,s
visibility timeout, six times the 60\,s Lambda timeout, so that a message is
never redelivered while it is still being processed. The pre-check accepts an
upload only if both images decode completely (\texttt{Image.open(...).load()}
without error); a truncated or corrupt file is rejected and retried until it
reaches the dead-letter queue.

\subsection{Alerting and Operational Response}\label{sec:alerting}

A dashboard is only useful if somebody is told when it turns red. Two alarms
are configured against the \texttt{vericloud-ops-alarm} topic
(Table~\ref{tab:alarms}); the notification that reaches the operations inbox is
shown in Fig.~\ref{fig:alarm}.

\begin{table}[htbp]
  \centering
  \small
  \caption{Operational alarms and their rationale.}
  \label{tab:alarms}
  \begin{tabular}{llll}
    \toprule
    Alarm & Condition & Period & Rationale \\
    \midrule
    Worker errors & Errors $>$ 5 & 5 min & Only ocr / face / notify \\
    DLQ depth & Visible $>$ 0 & 5 min & A package was rejected for good \\
    Workflow failures & Failed $>$ 0 & 5 min & The state machine aborted \\
    \bottomrule
  \end{tabular}
\end{table}

The three alarms are deliberately not redundant. Pre-check rejections are an
expected business outcome and are therefore observed through dead-letter depth
rather than through Lambda error counts; the error alarm covers only the workers
that should never fail (\texttt{vericloud-ocr}, \texttt{vericloud-face} and
\texttt{vericloud-notify}), so that a corrupted upload cannot make it ring.
Queue depth is the signal that actually matters for quarantine, because a message
reaches the dead-letter queue only after the configured number of retries has
been exhausted; the workflow alarm catches executions that abort before any
worker-level error is recorded.

\begin{figure}[htbp]
  \centering
  % TODO(E) 11-20: 上传 figures/fig_e_alarm.png 后取消注释
  % \includegraphics[width=0.85\linewidth]{figures/fig_e_alarm.png}
  \caption{Alarm notification delivered to the operations inbox.}
  \label{fig:alarm}
\end{figure}

\subsection{Results}\label{sec:results}

Table~\ref{tab:loadtest} summarises two consecutive runs of the same workload.
% TODO(E) 11-25: 用下面模板填，禁止先写结论
% Round 1 processed X/80 valid packages and isolated Y/20 corrupted uploads;
% round 2 processed X'/80 and isolated Y'/20. {No | N} executions terminated in
% an unhandled state. Median latency was ... s (p50) and ... s (p99).

\begin{table}[htbp]
  \centering
  \small
  \caption{Load test results (100 applications in 60\,s, two rounds).}
  \label{tab:loadtest}
  % TODO(E): 压测后取消注释——表格由 tests/collect_results.py 自动生成
  % \input{tables/tab_e_loadtest.tex}
\end{table}

% TODO(E): 两张图，图片上传后再取消注释
% \begin{figure}[htbp]
%   \centering
%   \includegraphics[width=0.95\linewidth]{figures/fig_e_dashboard.png}
%   \caption{CloudWatch dashboard during the load test.}
%   \label{fig:dashboard}
% \end{figure}
%
% \begin{figure}[htbp]
%   \centering
%   \includegraphics[width=0.8\linewidth]{figures/fig_e_latency.pdf}
%   \caption{End-to-end latency distribution for valid applications.}
%   \label{fig:latency}
% \end{figure}

\subsection{Discussion and Limitations}\label{sec:limits}

Three limitations qualify the results. First, the sample size is small: with one
hundred packages per round, the p99 estimate is determined by a single
observation and should be read as indicative rather than stable, which is why p50
and p95 are reported alongside it. Second, all measurements come from one
account in one region under the resource quotas of a teaching sandbox; a
production deployment would face different concurrency ceilings and a different
network path to Rekognition. Third, the corrupted-file scenarios are injected at
the object level, so the study does not exercise infrastructure-level faults such
as a regional outage of a managed service, and the corresponding recovery
behaviour is out of scope.

A further caveat concerns the fixtures themselves. Most documents are rendered
synthetically, which keeps personal data out of the repository but also removes
the lighting, tilt and sensor noise of a real capture. The ten normal packages
are drawn from a small number of enrolled identities
(\texttt{7} distinct synthetic identities), so the suite measures pipeline
behaviour rather than the demographic accuracy of the face model.
% TODO(E): 若最终换用组员真人照片，请把 7 改成实际人数

% TODO(E) 11-08 之后：若已打印并重拍 5 张 normal，替换为下面两句（否则不要写）：
% A small number of normal cases are printed at 300 dpi and re-photographed with
% a smartphone under ambient light and a 10--15 degree tilt, so that the reported
% OCR accuracy is not an artefact of digitally rendered card images.
